% Generated draft: requires mathematical and source review.
\documentclass[11pt]{article}
\usepackage[margin=25mm]{geometry}
\usepackage{fontspec}
\setmainfont{Latin Modern Roman}
\usepackage{amsmath,amssymb,mathtools,unicode-math}
\setmathfont{Latin Modern Math}
\usepackage{enumitem,needspace,xurl,hyperref}
\setlength{\parindent}{0pt}
\setlength{\parskip}{5pt}
\setlength{\emergencystretch}{3em}
\newcommand{\N}{\mathbb{N}}
\newcommand{\Z}{\mathbb{Z}}
\newcommand{\Q}{\mathbb{Q}}
\newcommand{\R}{\mathbb{R}}
\newcommand{\C}{\mathbb{C}}
\newcommand{\F}{\mathbb{F}}
\newcommand{\A}{\mathbb{A}}
\newcommand{\PP}{\mathbb{P}}
\newcommand{\E}{\mathbb{E}}
\newcommand{\eps}{\varepsilon}
\newcommand{\dd}{\,\mathrm{d}}
\begin{document}
\section*{UnsolvedMath ID: 30001770}
\textbf{Source problem code:} catalogue entry
\subsection{Definitions and mathematical statement}
Let \(k\) denote a coefficient field, whose characteristic and other hypotheses are not specified in the supplied record. For a finite group \(G\), write \(kG\) for its group algebra and \(B_0(kG)\) for its principal block, namely the block containing the trivial \(kG\)-module \(k\). A principal-block simple module is a nonzero simple \(kG\)-module \(S\) on which the central idempotent of \(B_0(kG)\) acts as the identity. Define group cohomology by
\[
H^i(G,S)=\operatorname{Ext}_{kG}^{i}(k,S),
\]
where \(k\) has the trivial \(G\)-action.

For a fixed coefficient field \(k\), the existence question in the supplied statement can be written as follows: does there exist a finite group \(G\) and a simple module \(S\) belonging to \(B_0(kG)\) such that
\[
\forall i\in\Z_{>0},\qquad H^i(G,S)=0?
\]
The same question is asked with \(G\) restricted to a symmetric group \(\operatorname{Sym}(\Omega)\), where \(\Omega\) is a finite set. The same module \(S\) must satisfy the vanishing condition simultaneously in every positive degree; no condition in degree zero is stated.

\textit{Unresolved qualifications.} The record does not specify the coefficient field, its characteristic, its splitting properties, or the intended quantifier over coefficient fields. It also supplies no hypothesis that the characteristic divides \(|G|\), no exclusion of the trivial module or trivial group, and no restriction on the degree of the symmetric group. Consequently, the displayed question exposes the missing coefficient parameter and does not determine a unique intended modular-representation conjecture. The phrase ``or be verified'' does not identify a separate mathematical assertion. The background names the Linnell conjecture but gives no precise formulation establishing its equivalence to this existence question.
\subsection{Short English statement}
Can a simple module in the principal block of a finite group have zero cohomology in every positive degree? The record also asks this for symmetric groups. The coefficient field and necessary restrictions are unspecified, so the intended conjectural formulation requires human review.
\subsection{Sources}
\begin{itemize}
\item Supplied local UnsolvedMath registry record, permanent ID \(30001770\), exact source problem code \texttt{catalogue entry}, title: \emph{Cohomology Vanishing for Principal-Block Simple Modules}. The record gives proposed year 2011 and identifies no individual proposer.
\item Supplied report citation: \emph{Mini-Workshop: Modular Representations of Symmetric Groups and Related Objects} (2011), Oberwolfach Reports, pp.~1177--1205, DOI \texttt{10.4171/owr/2011/21}. Source URL and literature-source link, labelled ``original OWR report'': \url{https://doi.org/10.4171/owr/2011/21}. The supplied background describes Oberwolfach Reports as documenting mathematical workshops and their open-problem discussions; it associates this question with the Linnell conjecture. No problem-specific page or quoted report passage is supplied.
\item Supplied literature assessment, checked on \textbf{2026-08-21}: the record labels the problem open, says that the report lists the Linnell conjecture as open, and reports finding no source establishing vanishing cohomology for all simple modules of finite or symmetric groups. These are supplied dated curation statements, not independently verified mathematical facts or a new literature assessment. The supplied background explicitly requests independent verification of time-sensitive status and cited claims.
\item The registry supplies the editorial label \texttt{corrected\_verified} and states that manual semantic, typographic, and LaTeX review confirmed a meaning-preserving cleanup. This is supplied editorial metadata, not independent verification of the mathematical formulation in this draft.
\end{itemize}
\end{document}
